% Generated draft: requires mathematical and source review.
\documentclass[11pt]{article}
\usepackage[margin=25mm]{geometry}
\usepackage{fontspec}
\setmainfont{Latin Modern Roman}
\usepackage{amsmath,amssymb,mathtools,unicode-math}
\setmathfont{Latin Modern Math}
\usepackage{enumitem,needspace,xurl,hyperref}
\setlength{\parindent}{0pt}
\setlength{\parskip}{5pt}
\setlength{\emergencystretch}{3em}
\newcommand{\N}{\mathbb{N}}
\newcommand{\Z}{\mathbb{Z}}
\newcommand{\Q}{\mathbb{Q}}
\newcommand{\R}{\mathbb{R}}
\newcommand{\C}{\mathbb{C}}
\newcommand{\F}{\mathbb{F}}
\newcommand{\A}{\mathbb{A}}
\newcommand{\PP}{\mathbb{P}}
\newcommand{\E}{\mathbb{E}}
\newcommand{\eps}{\varepsilon}
\newcommand{\dd}{\,\mathrm{d}}
\begin{document}
\section*{UnsolvedMath ID: 30003720}
\textbf{Source problem code:} catalogue entry
\subsection{Definitions and mathematical statement}
\textbf{Scope and unresolved definitions.} The supplied record concerns tropical principal component analysis (tropical PCA), tropical line segments in a space of trees, and possible connections between tropical geometry and deep learning. It does not specify the tree space, tropical arithmetic convention, metric, PCA objective, or meaning of robustness. Consequently, it specifies a research programme rather than a fully quantified conjecture.

For purposes of stating the requested tasks schematically, write \(\mathcal T\) for the intended tree space, \(D\) for an input collection of trees, and \(\operatorname{Top}(T)\) for the combinatorial topology of a tree \(T\in\mathcal T\). Write \(\operatorname{tseg}(T_0,T_1)\) for the intended tropical line segment with endpoints \(T_0,T_1\). These symbols are placeholders: the record does not determine their precise definitions, whether the segment remains in \(\mathcal T\), or whether topology retains rooting and leaf labels or permits unresolved vertices and zero-length edges.

\textbf{Requested mathematical investigation.}
\begin{enumerate}
\item For tropical PCA applied to \(D\), quantify how its fitted principal object and associated output change under outliers and perturbations. The input distance, output distance, outlier model, perturbation model, principal-object dimension, and treatment of nonunique fits must be specified before a sensitivity bound can be formulated. The record does not distinguish perturbations of observations from perturbations of the fitting procedure.
\item Describe the tree topologies occurring along \(\operatorname{tseg}(T_0,T_1)\) and how transitions between them occur. A precise formulation must specify the tree model, segment construction, and treatment of endpoints, degeneracies, and simultaneous transitions; these choices are absent from the record.
\item Determine whether tropical geometry yields meaningful mathematical connections to deep learning. The record supplies neither a class of network architectures nor a proposed correspondence, performance criterion, or theorem expressing what constitutes a meaningful or useful connection.
\end{enumerate}
No numerical robustness bound, particular topology classification, or deep-learning equivalence is asserted here. The supplied information does not determine such a statement or its quantifiers.
\subsection{Short English statement}
Investigate the sensitivity of tropical PCA and tropical tree segments to outliers and perturbations, describe topology changes along those segments, and seek precise connections with deep learning. The source leaves the mathematical models and success criteria unspecified, so this remains a statement-only research programme requiring clarification and human mathematical review.
\subsection{Sources}
\begin{itemize}
\item Local UnsolvedMath registry, supplied complete record: permanent ID \(30003720\), exact source problem code \texttt{catalogue entry}, title ``Robustness and Tree Topology in Tropical Principal Components.'' No individual proposer is recorded; the proposed year is recorded as 2018.
\item Supplied report citation: \emph{Statistics for Data with Geometric Structure} (2019), Oberwolfach Reports, DOI 10.4171/owr/2018/3, pp.~125--186. \url{https://doi.org/10.4171/owr/2018/3}. The supplied background describes Oberwolfach Reports as documenting mathematical workshops and their open-problem discussions. The citation year 2019 is retained alongside the recorded proposed year 2018 and the report identifier 3/2018; their relationship has not been independently checked.
\item Supplied literature source: ``Yoshida et al., tropical PCA.'' \url{https://academic.oup.com/bioinformatics/article/36/17/4590/5855129}.
\item Supplied literature source: ``Cox, tropical line-segment topologies.'' \url{https://arxiv.org/abs/2302.03611}.
\item Supplied literature source: ``Oberwolfach Report 3/2018.'' \url{https://doi.org/10.4171/owr/2018/3}.
\item Literature assessment dated exactly \texttt{2026-08-21}: the supplied curation labels the problem ``partially solved'' and reports substantial later results on tropical PCA algorithms and the behavior and classification of tree topologies along tropical line segments, while describing outlier sensitivity, perturbation theory, and links to deep learning as unresolved. This is supplied dated curation derived from the OWR material, not an independently verified literature assessment; time-sensitive status and cited claims require independent verification.
\item The registry labels the statement \texttt{corrected\_verified} and says that mathematical content was preserved while notation, typography, and sentence structure were repaired for publication. This is supplied editorial metadata, not independent mathematical verification of this draft.
\end{itemize}
\end{document}
